\documentclass{article}


% Use the official NeurIPS 2026 style file for submission. Omit the final/preprint option for anonymous review.
\usepackage{neurips_2026}

\usepackage{subcaption}
\usepackage[utf8]{inputenc} % allow utf-8 input
\usepackage[T1]{fontenc}    % use 8-bit T1 fonts
\usepackage{hyperref}       % hyperlinks
\usepackage{url}            % simple URL typesetting
\usepackage{booktabs}       % professional-quality tables
\usepackage{amsfonts}       % blackboard math symbols
\usepackage{nicefrac}       % compact symbols for 1/2, etc.
\usepackage{microtype}      % microtypography
\usepackage{xcolor}         % colors
\usepackage{amsmath}        % for \eqref and math environments
\usepackage{multirow}      % multi-row cells in tables
\usepackage{graphicx}
\usepackage{array}
\usepackage{longtable}

% \usepackage[nonatbib]{neurips_2024}
% \newcommand*\samethanks[1][\value{footnote}]{\footnotemark[#1]}
\title{VisText-VLM2Plastic: A Bilingual Benchmark for Visual Grounding and Environmental Reasoning in Vision-Language Models}

% \author{%
% author $^{2}$\thanks{Equal Contribution} \quad author$^{1,3*}$ \quad author$^{3}$ \\ 
% $^1$United International University \quad $^2$American International University Bangladesh \quad $^3$lAB NAME\\
% \texttt{m1}
% \texttt{m2}
% \texttt{m3}
% }


\begin{document}


\maketitle

% Avoid manual negative vertical spacing in the conference template.

\begin{abstract}
Vision-language models (VLMs) increasingly support open-ended visual reasoning, yet their behavior in low-resource languages and sustainability-oriented domains remains underexplored. We introduce a controlled English--Bangla plastic-waste VQA benchmark comprising 402 image records derived from 136 base images and 2,412 question--answer pairs across five waste categories and three task types: recognition, visual description, and environmental reasoning. We evaluate GPT-4.1 mini, Gemini 2.5 Flash, Qwen2.5-VL-7B-Instruct, and Gemma-3-4B zero-shot using task-specific answer-quality metrics, bilingual consistency, and image ablation. Gemini 2.5 Flash yields the highest recognition point estimate, with 93.5\% accuracy and 91.3\% macro-F1. Across models, English normalized performance exceeds Bangla by 3.78 points on average, with the largest observed gap for environmental reasoning. Removing the image reduces normalized performance by 66.29 points for recognition, 52.89 for description, and 11.42 for environmental reasoning. The smaller environmental-reasoning ablation gap indicates that, in this benchmark, many environmental questions remain partly answerable from linguistic and domain priors. These results show why answer quality alone is insufficient for multilingual multimodal evaluation and motivate reporting cross-lingual consistency and sensitivity to visual evidence alongside conventional metrics.
\end{abstract}

\section{Introduction}
\label{sec:introduction}

Plastic waste creates persistent environmental burdens through litter accumulation, drainage blockage, and movement into waterways. Computer vision has therefore been widely explored for waste classification and sorting, while recent vision-language models (VLMs) provide a more flexible alternative by jointly processing images and natural-language instructions. \citet{malla2025wasterecognition} showed that contemporary VLMs can achieve competitive zero-shot waste recognition and benefit substantially from prompt engineering.

Recognition alone, however, does not establish environmental understanding. A useful system should be able to identify the visible waste item, describe only visually supported properties, and connect the observed object to scientifically reasonable environmental concerns. These abilities are distinct: a model may correctly recognize an object while hallucinating unsupported visual details or producing generic environmental explanations that do not depend on the image.

Visual Question Answering (VQA) provides a natural framework for studying these distinctions \citep{antol2015vqa}. However, prior work has shown that high VQA accuracy can arise from language priors rather than genuine visual understanding \citep{goyal2017makingv}. This issue is particularly relevant for environmental questions, which may often be answered plausibly from pretrained knowledge without using the supplied image.

The problem becomes more important in multilingual settings. Bangla remains underrepresented in multimodal evaluation despite recent resources such as ChitroJera \citep{barua2024chitrojera}, BanglaProtha \citep{fahim2026banglaprotha}, and BanglaMedVQA \citep{ahmed2026banglamedvqa}. These studies show that strong general-purpose multimodal systems can degrade in specialized Bangla contexts.

We therefore develop a controlled English--Bangla benchmark for plastic-waste understanding organized into three task types that become progressively more open-ended:
\begin{equation}
\text{Recognition}
\;\longrightarrow\;
\text{Visual Description}
\;\longrightarrow\;
\text{Environmental Reasoning}.
\label{eq:reasoning_hierarchy}
\end{equation}

The benchmark addresses three research questions:
\begin{itemize}
    \item[\textbf{RQ1.}] How does VLM performance change across recognition, visual description, and environmental reasoning tasks?
    \item[\textbf{RQ2.}] How large is the performance gap between semantically aligned English and Bangla questions, and how does that gap vary across increasingly open-ended task types?
    \item[\textbf{RQ3.}] Within this benchmark, to what extent does performance on environmental questions decrease when visual input is removed?
\end{itemize}

Our contributions are:
\begin{itemize}
    \item We introduce a controlled English--Bangla diagnostic benchmark for plastic-waste understanding with paired questions for the same visual content and semantic intent.
    \item We evaluate proprietary and open VLMs across recognition, visual description, and environmental reasoning using classification, semantic, lexical, and cross-lingual metrics.
    \item We introduce an image-ablation protocol that measures how much benchmark performance changes when visual input is removed. We interpret this quantity as task-level sensitivity to visual evidence rather than as a standalone proof of causal visual grounding.
    \item We analyze bilingual consistency and qualitative failure modes, showing why multilingual environmental VQA is not adequately characterized by a single aggregate answer-quality score.
\end{itemize}

\section{Related Work}
\label{sec:related_work}

\subsection{Visual Question Answering and Language Priors}

VQA requires models to jointly reason over images and natural-language questions \citep{antol2015vqa}. However, \citet{goyal2017makingv} showed that models can exploit question regularities and language priors, producing strong scores without sufficiently interpreting the image. Our work follows this diagnostic perspective but directly removes the image to quantify modality dependence across task types.

\subsection{Bangla Multimodal Evaluation}

Bangla multimodal evaluation remains comparatively limited. ChitroJera introduced regionally relevant Bangla VQA data \citep{barua2024chitrojera}, BanglaProtha studied Bengali cultural understanding across VLMs \citep{fahim2026banglaprotha}, and BanglaMedVQA evaluated domain-specific medical reasoning in Bangla \citep{ahmed2026banglamedvqa}. Our benchmark differs by focusing on environmental reasoning and by using semantically aligned English--Bangla questions for the same images.

\subsection{Vision-Language Models for Waste Understanding}

Recent work has extended waste analysis beyond supervised image classification. \citet{malla2025wasterecognition} evaluated VLMs for zero-shot waste recognition and reported substantial gains from prompt optimization. WasteAssistant further introduced waste-oriented VQA and regulation-aware reasoning across multiple waste categories \citep{kataruka2026wasteassistant}. In contrast, our benchmark focuses on multilingual robustness and image sensitivity, asking whether Bangla performance differs from English and how strongly environmental-answer quality changes when visual input is removed.
\begin{table}[t] \centering \caption{Benchmark composition by plastic-waste category.} \label{tab:benchmark_composition} \begin{tabular}{lrr} \toprule \textbf{Category} & \textbf{Image Records} & \textbf{Fraction} \\ \midrule Plastic wrapper / flexible packaging & 224 & 55.7\% \\ Plastic bottle & 70 & 17.4\% \\ Plastic cup & 48 & 11.9\% \\ Plastic straw & 36 & 9.0\% \\ Plastic bag / plastic sheet & 24 & 6.0\% \\ \midrule \textbf{Total} & \textbf{402} & \textbf{100\%} \\ \bottomrule \end{tabular} \end{table}
\section{Dataset Benchmark}
\label{sec:benchmark}
\begin{figure}[t]
    \centering
    \includegraphics[width=\linewidth]{figures/benchmark_pipeline.png}
    \caption{Benchmark and evaluation pipeline. Each grouped plastic-waste image produces aligned English and Bangla recognition, description, and environmental-reasoning questions. Four VLMs are evaluated under image+question and question-only conditions, followed by multilingual, grounding, and error analyses.}
    \label{fig:benchmark_pipeline}
\end{figure}
The benchmark contains 402 image records derived from 136 underlying base images. Multiple records correspond to augmented variants of the same source image; therefore, a \texttt{source\_base\_id} is preserved for every record, and all augmentations originating from the same source image are treated as a single group during dataset splitting. The final benchmark contains five dominant plastic-waste categories, whose record-level distribution is summarized in Table~\ref{tab:benchmark_composition}. The dataset is imbalanced, with flexible wrappers comprising more than half of the image records, motivating the use of macro-averaged metrics in addition to micro accuracy. The dataset is split into 282/58/62 train, validation, and test image records, corresponding to 95/20/21 base images and 1,692/348/372 QA pairs, respectively, and no \texttt{source\_base\_id} appears in more than one partition. This grouping prevents visually near-identical augmented samples from appearing across development and evaluation sets, which could otherwise inflate performance. Because the 62 test records correspond to only 21 independent base images, uncertainty is estimated at the base-image-group level rather than by treating augmented records as independent. The complete benchmark and evaluation flow is summarized in Figure~\ref{fig:benchmark_pipeline}.

% TODO(author): Add per-category counts at the independent base-image level, especially for the 21-base-image test set. Record-level counts alone do not show effective class support.

Each image record is associated with six QA pairs: English recognition, Bangla recognition, English visual description, Bangla visual description, English environmental reasoning, and Bangla environmental reasoning. This produces 2,412 QA pairs, evenly divided into 1,206 English and 1,206 Bangla instances. Question construction follows a controlled template-based procedure in which surface wording varies while task semantics are preserved. All 2,412 questions have unique surface strings in the current release. Reference answers are generated from a constrained category-conditioned answer bank containing 100 distinct answer strings across the two languages. The design therefore favors controlled paired comparison over naturalistic linguistic diversity. Repeated answer semantics can also make some questions, especially generic environmental questions, partly recoverable from language and domain priors; we account for this limitation when interpreting the image-ablation results. The overall data-collection and construction workflow is illustrated in Appendix~\ref{appendix:dataset}, Figure~\ref{fig:data_collection_pipeline}, and the trade-offs are discussed in Appendix~\ref{sec:limitations}.

% \begin{table}[t]
% \centering
% \caption{Dataset split statistics.}
% \label{tab:dataset_split}
% \begin{tabular}{lrrr}
% \toprule
% \textbf{Split} & \textbf{Image Records} & \textbf{Base Images} & \textbf{QA Pairs} \\
% \midrule
% Train      & 282 & 95  & 1,692 \\
% Validation & 58  & 20  & 348 \\
% Test       & 62  & 21  & 372 \\
% \midrule
% \textbf{Total} & \textbf{402} & \textbf{136} & \textbf{2,412} \\
% \bottomrule
% \end{tabular}
% \end{table}

\begin{table}[t]
\centering
\caption{Zero-shot recognition performance.}
\label{tab:recognition_results}
\begin{tabular}{lcc}
\toprule
\textbf{Model} & \textbf{Accuracy (\%)} & \textbf{Macro-F1 (\%)} \\
\midrule
GPT-4.1 mini
& 91.9
& 89.2 \\

Gemini 2.5 Flash
& \textbf{93.5}
& \textbf{91.3} \\

Qwen2.5-VL-7B
& 88.7
& 85.6 \\

Gemma 3 4B
& 85.5
& 81.4 \\
\bottomrule
\end{tabular}
\end{table}
\section{Experimental Setup}
\label{sec:experimental_setup}
Experiments were conducted in Google Colab using an NVIDIA Tesla T4 GPU, with the pipeline implemented in Python using the required deep-learning and vision-language libraries. The T4 GPU was used for dataset preprocessing, local inference with the open-source VLMs, and metric computation, while the proprietary models were accessed through their respective APIs. We evaluate four representative VLMs spanning proprietary and openly available model families: GPT-4.1 mini, using the fixed API snapshot \texttt{gpt-4.1-mini-2025-04-14} where available, which supports image inputs and improved multimodal understanding and instruction following \citep{openai2025gpt41}; Gemini 2.5 Flash, which serves as a proprietary multimodal baseline optimized for efficient reasoning \citep{google2025gemini25flash}; Qwen2.5-VL-7B-Instruct, which supports dynamic-resolution visual processing and multilingual vision-language interaction \citep{bai2025qwen25vl}; and the multimodal Gemma-3-4B Instruction-Tuned model, which provides vision understanding and expanded multilingual capabilities \citep{gemmateam2025gemma3}. The comparison is intended to represent different deployment regimes, model scales, and accessibility conditions rather than to establish a universal ranking, since model families and API implementations may evolve over time. Additional details regarding the experimental setup and evaluation procedures are provided in Appendix~\ref{appendix:metrics}.

% TODO(author): Report exact revision/checkpoint identifiers for Qwen2.5-VL-7B-Instruct and Gemma-3-4B, and the exact Gemini model/API version and inference date if available.


\section{Results}
\label{sec:results}


\subsection{Zero-Shot Recognition}
\label{subsec:recognition_results}

Table~\ref{tab:recognition_results} reports the zero-shot recognition performance of all four evaluated models.
Recognition yields high category accuracy across all four evaluated models. Gemini 2.5 Flash has the highest accuracy point estimate at 93.5\%, followed by GPT-4.1 mini at 91.9\%. Given the small number of independent test base images, these point estimates should not be interpreted as a robust model ranking without paired uncertainty on the differences. The difference between accuracy and macro-F1 is consistent with the class imbalance described in Section~\ref{sec:benchmark}: performance on the dominant flexible-wrapper class has greater influence on micro accuracy.

Category-level results complement the aggregate metrics and reveal greater sensitivity for bag and straw examples, which have fewer samples and can visually overlap with flexible packaging. Figure~\ref{fig:category_performance} presents the category distribution together with per-category recognition performance.
\begin{table*}[t]
\centering
\caption{Open-ended response quality across description and environmental-reasoning tasks. }
\label{tab:generation_results}
\begin{tabular}{lcccccc}
\toprule
&
\multicolumn{3}{c}{\textbf{Description}}
&
\multicolumn{3}{c}{\textbf{Environmental Reasoning}}
\\
\cmidrule(lr){2-4}
\cmidrule(lr){5-7}
\textbf{Model}
&
\textbf{BERT-F1}
&
\textbf{BLEU-4}
&
\textbf{ROUGE-L}
&
\textbf{BERT-F1}
&
\textbf{BLEU-4}
&
\textbf{ROUGE-L}
\\
\midrule

GPT-4.1 mini
& 0.868
& 0.431
& 0.619
& 0.840
& 0.347
& 0.566 \\

Gemini 2.5 Flash
& \textbf{0.881}
& \textbf{0.452}
& \textbf{0.641}
& \textbf{0.855}
& \textbf{0.369}
& \textbf{0.588} \\

Qwen2.5-VL-7B
& 0.834
& 0.398
& 0.587
& 0.802
& 0.311
& 0.531 \\

Gemma 3 4B
& 0.804
& 0.361
& 0.552
& 0.769
& 0.276
& 0.498 \\

\bottomrule
\end{tabular}
\end{table*}
\begin{figure}[t]
    \centering
    \includegraphics[width=\linewidth]{figures/category_performance.png}
    \caption{Category distribution and per-category recognition performance. Error bars represent 95\% cluster-bootstrap confidence intervals over source base images.}
    \label{fig:category_performance}
\end{figure}

\subsection{Description and Environmental Generation}
\label{subsec:generation_results}
Table~\ref{tab:generation_results} summarizes BERTScore-F1, BLEU-4, and ROUGE-L for the description and environmental-reasoning tasks. Across all three automatic metrics, visual-description responses obtain higher answer--reference similarity than environmental-reasoning responses, and the two proprietary systems have higher point estimates than the selected 7B and 4B open checkpoints. These differences should not be interpreted as isolating ``reasoning difficulty,'' because the task types also differ in answer length, answer space, linguistic generation demands, and required domain knowledge. BERTScore values are substantially higher than BLEU-4 scores. This is expected rather than contradictory. A model may describe a wrapper as ``thin flexible packaging'' when the reference uses ``lightweight plastic wrapper with a sheet-like form.'' The two responses preserve similar semantic content while sharing relatively few four-grams. This distinction motivates using BERTScore as the primary open-ended similarity measure \citep{zhang2020bertscore}, with BLEU-4 and ROUGE-L serving as complementary lexical metrics \citep{papineni2002bleu,lin2004rouge}.

The lower environmental-reasoning similarity scores are nevertheless informative because these answers must combine object-related context with domain knowledge while avoiding unsupported toxicological, chemical, or causal claims. We therefore treat the comparison as a difference between task types, not as a controlled estimate of reasoning complexity.
\begin{figure}[t]
    \centering
    \includegraphics[width=\linewidth]{figures/language_gap.png}
    \caption{English--Bangla performance gap across recognition, description, and environmental reasoning. Error bars represent 95\% cluster-bootstrap confidence intervals over source base images. The task types are categorical; connecting lines are used only as a visual guide.}
    \label{fig:language_gap}
\end{figure}
% TODO(author): Consider regenerating Figure~\ref{fig:language_gap} as grouped points or bars rather than a connected line plot, because the task types are categorical.

\begin{figure}[t]
    \centering
    \includegraphics[width=\linewidth]{figures/grounding_ablation.png}
    \caption{Visual-ablation analysis. Recognition and description exhibit substantially larger multimodal gains than environmental reasoning in the experimental results.}
    \label{fig:grounding_ablation}
\end{figure}
\subsection{English--Bangla Performance Gap}
\label{subsec:language_gap}
We next examine whether semantically aligned English and Bangla questions produce equivalent model behavior. To place the task-specific primary metrics on a common 0--100 reporting scale, we define
\begin{equation}
S_{\mathrm{norm}}(t)=
\begin{cases}
100\,\operatorname{Acc}, & t=\text{recognition},\\
100\,\operatorname{BERTScoreF1}, & t\in\{\text{description},\text{environmental}\}.
\end{cases}
\label{eq:normalized_primary_score}
\end{equation}
When a single model--language summary is reported, we take the unweighted mean of $S_{\mathrm{norm}}$ across the three task types. BLEU-4 and ROUGE-L are reported separately and are not included in this normalized summary. Model-level normalized performance and the corresponding English--Bangla gap are reported in Table~\ref{tab:language_model_gap}.
% TODO(author): Confirm that Equation~\ref{eq:normalized_primary_score} exactly matches the evaluation script used to produce Tables~\ref{tab:language_model_gap}, \ref{tab:language_task_gap}, and \ref{tab:visual_ablation}.
\begin{table}[t]
\centering
\caption{Language-level normalized primary performance on a 0--100 scale; model-level values are unweighted means across the three task types.}
\label{tab:language_model_gap}
\begin{tabular}{lccc}
\toprule
\textbf{Model}
& \textbf{English}
& \textbf{Bangla}
& \textbf{EN--BN Gap} \\
\midrule
GPT-4.1 mini
& 89.23
& 85.87
& 3.36 \\

Gemini 2.5 Flash
& \textbf{90.60}
& \textbf{87.43}
& \textbf{3.17} \\

Qwen2.5-VL-7B
& 85.90
& 82.23
& 3.67 \\

Gemma 3 4B
& 83.10
& 78.20
& 4.90 \\

\midrule
\textbf{Mean}
& --
& --
& \textbf{3.78} \\
\bottomrule
\end{tabular}
\end{table}
All four evaluated models exhibit an English advantage in the normalized primary score. The observed gap is smallest for Gemini 2.5 Flash and largest for Gemma 3 4B. More informative than the overall difference is how the language gap varies by task type; Table~\ref{tab:language_task_gap} reports this comparison.
\begin{table}[t]
\centering
\caption{Language gap by task type using the normalized primary score.}
\label{tab:language_task_gap}
\begin{tabular}{lccc}
\toprule
\textbf{Task}
& \textbf{English}
& \textbf{Bangla}
& \textbf{Gap} \\
\midrule
Recognition
& 91.53
& 87.90
& 3.63 \\

Description
& 86.43
& 82.88
& 3.55 \\

Environmental reasoning
& 83.68
& 79.53
& \textbf{4.15} \\
\bottomrule
\end{tabular}
\end{table}

Environmental reasoning exhibits the largest observed language gap. This pattern is consistent with greater cross-lingual divergence in the most open-ended task, but the difference should not be attributed solely to reasoning complexity because answer length, answer space, and generation requirements also vary across task types. Language-specific open-ended scores are reported in Table~\ref{tab:language_generation_metrics}; lexical metrics require additional caution.

\begin{table}[t]
\centering
\caption{Language-specific open-ended evaluation. }
\label{tab:language_generation_metrics}
\begin{tabular}{lccc}
\toprule
\textbf{Language}
& \textbf{BERT-F1}
& \textbf{BLEU-4}
& \textbf{ROUGE-L} \\
\midrule
English
& 0.842
& 0.382
& 0.574 \\

Bangla
& 0.806
& 0.329
& 0.531 \\
\bottomrule
\end{tabular}
\end{table}

The BLEU gap is proportionally larger than the BERTScore gap. This suggests that some Bangla responses may remain semantically appropriate despite greater surface-form variation. BLEU should therefore not be used alone to infer weaker Bangla reasoning. Figure~\ref{fig:language_gap} visualizes the English--Bangla performance gap across the three task types.




\subsection{Does the Image Actually Matter?}
\label{subsec:image_ablation_results}
Table~\ref{tab:visual_ablation} summarizes the central diagnostic experiment.
\begin{table}[t]
\centering
\caption{Visual-ablation analysis using the normalized primary score. VGG is the Image+Q minus Q-Only difference.}
\label{tab:visual_ablation}
\begin{tabular}{lccc}
\toprule
\textbf{Task}
&
\textbf{Image + Q}
&
\textbf{Q Only}
&
\textbf{VGG} \\
\midrule
Recognition
& 89.71
& 23.42
& \textbf{+66.29} \\

Description
& 84.65
& 31.76
& \textbf{+52.89} \\

Environmental reasoning
& 81.60
& 70.18
& \textbf{+11.42} \\
\bottomrule
\end{tabular}
\end{table}
Recognition behaves as expected: removing the image produces a large performance collapse. Description shows a similarly large reduction, indicating that reference-aligned answers to these benchmark questions are strongly sensitive to the presence of visual evidence.

Environmental reasoning behaves differently. Even without an image, the observed results retain 70.18 normalized performance compared with 81.60 under multimodal input, yielding a VGG of 11.42 points. This benchmark-specific finding indicates that the current environmental-question templates and reference semantics permit many broadly acceptable answers to be recovered from linguistic and domain priors. Terms such as ``plastic waste,'' ``improper disposal,'' and ``environmental problem'' can trigger plausible statements about litter, drainage systems, waterways, or wildlife even when the image is absent. Consequently, a high environmental-reasoning score does not, by itself, establish object-specific visual grounding. Conversely, a small VGG should not be interpreted as general evidence that a model ignores images; it may also reflect benchmark questions that are insufficiently discriminative without visual input.

This behavior is conceptually related to the language-prior problem identified in prior VQA research \citep{goyal2017makingv}. The issue is not necessarily that the generated environmental answer is false; rather, the \emph{source of evidence} supporting that answer may differ from what a multimodal benchmark is intended to measure. We therefore interpret VGG as a benchmark-conditioned sensitivity diagnostic rather than a standalone causal measure of visual grounding. The task-wise multimodal and question-only comparison is visualized in Figure~\ref{fig:grounding_ablation}.



\subsection{Cross-Lingual Consistency}
\label{subsec:bcs_results}
A model may achieve respectable aggregate scores in both languages while producing contradictory answers to paired questions. We therefore evaluate cross-lingual response consistency, with model-wise BCS results reported in Table~\ref{tab:bcs_results}.

The observed pattern indicates a monotonic decline in consistency across increasingly open-ended task types. Recognition is highly stable because outputs often reduce to the same short waste category. Description introduces greater linguistic variation, while environmental reasoning provides more opportunities for differences in specificity, factual content, and grounding. Because task format and response length change simultaneously, this trend should not be treated as an isolated effect of reasoning complexity. This analysis captures a failure mode that independent English and Bangla evaluation can conceal. For example, a model may identify an object as a plastic wrapper in English but use a broader or different waste category in Bangla while still generating environmentally plausible follow-up text. Both responses may obtain reasonable independent semantic scores while remaining mutually inconsistent.
\begin{table}[htb]
\centering
\caption{Bilingual Consistency Score (BCS). }
\label{tab:bcs_results}
\begin{tabular}{lccc}
\toprule
\textbf{Model}
&
\textbf{Rec.}
&
\textbf{Desc.}
&
\textbf{Env.} \\
\midrule
GPT-4.1 mini
& 0.94
& 0.89
& 0.84 \\

Gemini 2.5 Flash
& \textbf{0.96}
& \textbf{0.91}
& \textbf{0.87} \\

Qwen2.5-VL-7B
& 0.91
& 0.85
& 0.79 \\

Gemma 3 4B
& 0.87
& 0.81
& 0.73 \\
\bottomrule
\end{tabular}
\end{table}


%
\section{Discussion}
\label{sec:discussion}
%

\subsection{Task Type and Multilingual Robustness}
\label{subsec:task_multilingual_discussion}

The benchmark is organized across three progressively more open-ended task types. Recognition yields high closed-form category accuracy, whereas description and environmental reasoning receive lower open-ended answer--reference similarity scores. Because the tasks differ in answer space, response length, evaluation metric, and required domain knowledge, these differences should not be interpreted as a controlled measure of reasoning difficulty. Recognition accuracy is nevertheless insufficient by itself to characterize environmental multimodal capability, particularly for applications such as environmental education, citizen reporting, and waste-assistance systems.

The results also show that multilingual degradation varies by task type. Short recognition answers require comparatively little generation, whereas environmental questions require longer domain-specific responses. The English--Bangla gap is largest for environmental reasoning, suggesting that cross-lingual divergence is more visible in this open-ended setting. However, this comparison also reflects differences in response format and linguistic generation demands, so the gap should not be attributed solely to deeper reasoning. This observation is consistent with broader findings from low-resource multimodal evaluation; for example, BanglaMedVQA reports that specialized visual reasoning remains challenging even for strong general-purpose LVLMs \citep{ahmed2026banglamedvqa}. A methodological advantage of our benchmark is that English and Bangla questions are paired for the same image and semantic intent, reducing visual-content confounds in cross-lingual comparison.

\subsection{Visual Grounding, Metric Interpretation, and Model Comparison}
\label{subsec:grounding_metric_models}

The image-ablation experiment exposes a benchmark-design issue as well as a model-behavior issue. If a model can answer an environmental question nearly as well without seeing the image, the resulting score partly reflects environmental language knowledge and question priors rather than object-specific visual evidence. General environmental knowledge is useful, but it should not be interpreted as evidence that the supplied image was necessary for the answer. Accordingly, VGG measures performance sensitivity to image removal within this benchmark; it is not a universal measure of a model's grounding ability. Future benchmark extensions could reduce language-only shortcuts through counterfactual visual pairs, where similar question structures are paired with different waste categories and reference answers require category-specific consequences.

This grounding issue also affects the interpretation of automatic metrics. BLEU rewards local $n$-gram precision \citep{papineni2002bleu}, ROUGE-L measures longest-common-subsequence overlap \citep{lin2004rouge}, and BERTScore is more tolerant of semantically correct paraphrases \citep{zhang2020bertscore}. These metrics therefore capture different aspects of response quality, especially in bilingual evaluation where Bangla morphology and phrase variation can reduce lexical overlap. None of them directly establishes visual grounding or scientific correctness, so we treat them as complementary and combine them with modality ablation and a qualitative human audit.

The experimental point estimates place Gemini 2.5 Flash and GPT-4.1 mini above the selected open checkpoints on the reported metrics. However, this comparison spans different model sizes, training corpora, architectures, inference infrastructures, and deployment constraints, and should not be interpreted as evidence that proprietary VLMs are inherently superior. Open models remain advantageous for reproducibility, local deployment, privacy, modification, and scientific inspection. Future comparisons should include larger open checkpoints and more closely matched compute and latency budgets.

% 


% % 
% \section{Reproducibility}
% \label{sec:reproducibility}
% % 

% The evaluation pipeline exports raw predictions before metric computation. For each prediction, the released result file should contain the model identifier, base-image identifier, image filename, split, category, language, question type, question, reference answer, generated response, and inference condition.

% The accompanying code release should include:

% \begin{itemize}
%     \item dataset loading and grouping by \texttt{source\_base\_id};
%     \item zero-shot prompts and inference configurations;
%     \item exact recognition-normalization rules;
%     \item BERTScore, BLEU-4, ROUGE-L, and LaBSE implementations;
%     \item image-ablation inference;
%     \item source-group cluster-bootstrap confidence intervals;
%     \item bilingual-consistency analysis;
%     \item category-level evaluation; and
%     \item scripts for regenerating all tables and figures.
% \end{itemize}

% For API-based models, raw responses should be cached to reduce instability caused by repeated calls or subsequent model updates. For local models, package versions, checkpoint revisions, quantization settings, GPU type, inference precision, and decoding parameters should be recorded.


% % 
\section{Conclusion}
\label{sec:conclusion}
% 

We introduce a controlled bilingual framework for evaluating plastic-waste understanding in English and Bangla. The benchmark contains 402 image records derived from 136 base images and 2,412 paired QA instances spanning recognition, visual description, and environmental reasoning. Paired English--Bangla questions enable language comparisons on matched visual content, BCS measures cross-lingual response consistency, and image ablation measures how benchmark performance changes when visual input is removed.

The experimental results show high recognition accuracy, lower answer--reference similarity on the more open-ended tasks, an English--Bangla performance gap, and a much smaller image-ablation gain for environmental reasoning than for recognition or visual description. The last finding is benchmark-specific: it indicates that many current environmental questions can be answered plausibly from linguistic and domain priors, not that VLMs generally fail to use visual information.

For environmental multimodal evaluation, plausible answer generation should therefore be separated from evidence use. Future extensions should increase the number and diversity of independent source images, validate Bangla and category labels with human annotators, and introduce counterfactual or category-specific questions that cannot be answered reliably without the image.
\bibliographystyle{plainnat}
\bibliography{refs}
\clearpage
\appendix
\section{Extra Dataset Information}
\label{appendix:dataset}
The overall dataset collection and construction process is illustrated in Figure \ref{fig:data_collection_pipeline}
\begin{figure*}[t]
    \centering
    \includegraphics[width=\textwidth]{figures/data_collection_vqa_pipeline.pdf}
    \caption{Data collection and bilingual VQA benchmark construction pipeline. Plastic-waste images are collected from different locations in Dhaka, annotated using bounding boxes, organized into an annotated dataset, and converted into English--Bangla VQA pairs spanning recognition, visual description, and environmental reasoning.}
    \label{fig:data_collection_pipeline}
\end{figure*}
\subsection{Quality Control}
\label{subsec:quality_control}

The benchmark records image-quality and generation-status metadata. Of the 402 images, 385 are labeled \texttt{clear} and 17 \texttt{moderate}; none are retained as fully ambiguous. Seventeen images carry a low-image-quality warning, while one source has an empty label file and one has a missing label file.
% TODO(author): State explicitly how the empty and missing label-file cases were resolved and whether either case entered the final benchmark/evaluation set.

Our construction is intentionally conservative but semi-automated. Category mappings and QA generation are not equivalent to independent expert annotation of every image. This distinction is essential: the benchmark should be interpreted as a controlled diagnostic dataset rather than a naturalistic human-authored VQA corpus.
\subsection{Task Types}
\label{subsec:task_types}

The benchmark contains exactly 804 QA pairs for each task type.

\paragraph{Recognition (Closed-form).}
Questions ask which plastic-waste category is visually dominant. Answers are intentionally concise and category-centered.

\paragraph{Visual Description (Open-ended visual).}
Questions request observable properties such as general shape, flexibility, visible structure, openness, or deformation. Answers are designed to avoid claims that cannot be reliably inferred from visual evidence.

\paragraph{Environmental Reasoning (Open-ended domain).}
Questions ask why improper disposal may be environmentally problematic or what environmental concern is associated with the observed waste item. Reference answers use conservative concepts such as persistent litter, land and water pollution, movement into drains or waterways, wildlife exposure, and waste-management challenges.

Importantly, environmental answers do not claim exact degradation times, chemical toxicity, carbon footprint, contamination status, or direct medical consequences unless such evidence is externally established. 

\section{Additional Experimental Information}
\label{appendix:metrics}
\subsection{Zero-Shot Evaluation and Metrics}

All models are evaluated in a zero-shot setting without fine-tuning on the training or validation split. For the primary multimodal condition, each model receives the image together with the question in its original language and a short instruction encouraging direct, visually grounded, and scientifically conservative responses. A representative prompt is: ``Answer the question using the image. Use only visually supported information for recognition and description. For environmental questions, provide a concise scientifically conservative answer. Do not invent brands, exact chemical properties, locations, toxicity, or hidden objects.'' Temperature is set to 0 where supported, output length is capped to avoid unnecessary elaboration, and Bangla questions are processed directly without translation into English. For recognition, model outputs are normalized to the five target categories using a deterministic synonym mapping; for example, ``bottle,'' ``plastic bottle,'' and ``bottle-like plastic container'' are mapped to \texttt{plastic\_bottle}. Recognition performance is measured using micro accuracy,
\begin{equation}
\operatorname{Acc}
=
\frac{1}{N}
\sum_{i=1}^{N}
\mathbb{1}
\left[
\hat{y}_i = y_i
\right],
\label{eq:accuracy}
\end{equation}
where $y_i$ and $\hat{y}_i$ denote the reference and predicted categories, respectively, and macro-F1,
\begin{equation}
F1_{\mathrm{macro}}
=
\frac{1}{C}
\sum_{c=1}^{C}
F1_c,
\label{eq:macro_f1}
\end{equation}
where $C$ is the number of waste categories. Macro-F1 is reported because the benchmark is class-imbalanced and flexible-wrapper examples are dominant. For open-ended visual-description and environmental-reasoning responses, exact match is insufficient because semantically valid answers may differ in wording. We therefore combine semantic and lexical metrics. BERTScore-F1 is used as the primary automatic metric because it compares candidate and reference answers using contextual embeddings and is more tolerant of paraphrasing \citep{zhang2020bertscore}. BLEU-4 measures modified $n$-gram precision \citep{papineni2002bleu}, while ROUGE-L captures longest-common-subsequence overlap \citep{lin2004rouge}; both are reported as complementary lexical measures. Because English and Bangla differ in morphology, segmentation, and surface realization, lexical scores are interpreted cautiously and are not treated as perfectly language-neutral. The BERTScore encoder and scoring configuration can also affect cross-lingual comparability and must therefore be reported explicitly. For reproducibility, the experiments record the exact model or API snapshot, inference date, decoding parameters, package versions, and model configuration.
% TODO(author): Specify the exact BERTScore encoder/checkpoint, bert-score package version, IDF setting, rescale-with-baseline setting, and any language/tokenization configuration used for English and Bangla.
\subsection{Bilingual Consistency}
\label{subsec:bilingual_consistency}

Independent English and Bangla scores do not establish whether a model gives semantically equivalent answers to equivalent questions. We therefore compute a Bilingual Consistency Score (BCS). For recognition, consistency is defined as category agreement across paired English and Bangla outputs. For open-ended answers, paired responses are represented in a multilingual sentence-embedding space using LaBSE, which was developed to provide language-agnostic sentence representations across more than 100 languages \citep{feng2022labse}. Given paired English and Bangla outputs $e_i$ and $b_i$, respectively, we compute

\begin{equation}
s_i =
\cos
\left(
\operatorname{LaBSE}(e_i),
\operatorname{LaBSE}(b_i)
\right).
\label{eq:labse_similarity}
\end{equation}

After validating a semantic threshold $\tau$ on a manually reviewed validation subset, BCS is defined as

\begin{equation}
BCS
=
\frac{1}{N}
\sum_{i=1}^{N}
\mathbb{1}
\left[
s_i \geq \tau
\right].
\label{eq:bcs}
\end{equation}

The threshold $\tau$ is selected on a manually reviewed validation subset and then held fixed during test evaluation.
% TODO(author): Report the numerical value of tau, the size of the reviewed subset, annotator count, threshold-selection criterion, and (if available) a sensitivity analysis showing how BCS changes around the selected threshold.

\subsection{Visual-Grounding Ablation}
\label{subsec:grounding_ablation}

Our central diagnostic removes the image while leaving the question unchanged.

The multimodal condition is

\begin{equation}
\mathcal{M}:
(I,Q)
\rightarrow
A,
\label{eq:multimodal_condition}
\end{equation}

whereas the language-only condition is

\begin{equation}
\mathcal{L}:
Q
\rightarrow
A.
\label{eq:language_condition}
\end{equation}

For an evaluation metric $S$, we define Visual Grounding Gain (VGG) as

\begin{equation}
VGG
=
S(\mathcal{M})
-
S(\mathcal{L}).
\label{eq:vgg}
\end{equation}

Within a fixed benchmark and metric, a larger positive value indicates greater performance sensitivity to image removal. A small value can arise because the model relies weakly on the image, because the question can be answered from linguistic or domain priors, or both. VGG is therefore a diagnostic of benchmark-conditioned image sensitivity rather than a standalone causal measure of visual grounding.

This distinction is particularly important for environmental reasoning because general facts concerning plastic pollution may remain recoverable even when the model does not know whether the image contains a bottle, wrapper, cup, bag, or straw.


\subsection{Statistical Analysis}
\label{subsec:statistics}

Augmented variants originating from the same base image are statistically correlated. Treating all 62 test records as independent observations would therefore underestimate uncertainty. We use a cluster-bootstrap procedure based on \texttt{source\_base\_id}, rather than an ordinary QA-level bootstrap.
For uncertainty analyses, base-image groups are sampled with replacement using bootstrap resampling. All augmented records and associated QA pairs belonging to a sampled base image are retained together. Figures report percentile-based 95\% confidence intervals. Pairwise model comparisons use the same bootstrap samples across models, producing paired bootstrap estimates of performance differences. Small differences in point estimates should not be treated as robust model ordering unless the paired uncertainty on the difference supports that interpretation.
% TODO(author): If model-to-model differences are discussed as substantive, report paired 95% confidence intervals for those differences in a table or appendix.
\section{Additional Results}
\subsection{Error Analysis}
\label{subsec:error_analysis}

We manually audited a stratified subset of predictions covering models, languages, and question types. The annotation protocol distinguishes five failure categories. Because this analysis is conducted on a subset rather than the full test set, the resulting percentages are interpreted descriptively rather than as population estimates.
% TODO(author): Report the audited sample size, stratification/sampling rule, annotator count, whether failure categories can overlap, and any adjudication or agreement procedure.

\paragraph{Object misclassification.}
The dominant plastic item is assigned to an incorrect category.

\paragraph{Visual hallucination.}
The model invents a visible property, object, label, color, cap, opening, or condition that is not sufficiently supported by the image.

\paragraph{Generic environmental reasoning.}
The answer is environmentally plausible but nearly independent of the specific image or waste category.

\paragraph{Environmental hallucination.}
The model makes unsupported claims concerning toxicity, chemical composition, health consequences, or other scientific properties.

\paragraph{Bangla interpretation error.}
The model misinterprets the intent of the Bangla question or responds with semantically inappropriate content.

The observed qualitative error distribution is shown in Table~\ref{tab:error_distribution}.

\begin{table}[t]
\centering
\caption{Qualitative error distribution within the manually audited subset.}
\label{tab:error_distribution}
\begin{tabular}{lcc}
\toprule
\textbf{Error Type}
& \textbf{English}
& \textbf{Bangla} \\
\midrule
Object misclassification
& 6.5\%
& 9.7\% \\

Visual hallucination
& 5.4\%
& 8.1\% \\

Generic environmental reasoning
& 15.1\%
& \textbf{21.5\%} \\

Environmental hallucination
& 4.3\%
& 6.5\% \\

Bangla interpretation error
& --
& 7.5\% \\
\bottomrule
\end{tabular}
\end{table}

The observed error distribution identifies generic environmental reasoning as the dominant failure mode, particularly for Bangla. This behavior aligns with the visual-ablation analysis: models may possess broadly correct knowledge about plastic pollution while failing to condition that knowledge strongly on the precise visual evidence.

% The previous draft defined an additional H=C+R+G human score but did not report the resulting scores, sample size, or inter-annotator agreement. It is omitted here to avoid presenting an incomplete evaluation protocol. Reintroduce it only if the corresponding empirical results are available.

\section{Limitations and Threats to Validity}
\label{sec:limitations}
% 

\paragraph{Dataset size, diversity, and statistical power.}
The benchmark contains 136 underlying base images, and the test split contains only 21 independent base images. The 402 image records include augmented variants, so the effective visual diversity and independent test support are substantially smaller than the raw record count suggests. Cluster bootstrap addresses dependence among augmentations, but it cannot create additional independent visual diversity. The benchmark should therefore be interpreted as a focused diagnostic resource rather than a comprehensive representation of global plastic pollution, and small model-to-model differences should be treated cautiously.

\paragraph{Category imbalance and category overlap.}
Flexible wrappers account for 55.7\% of image records, whereas plastic bags represent only 6.0\%. Micro accuracy can therefore overstate performance on the dominant class. Macro-F1 and per-category evaluation reduce, but do not eliminate, this limitation. In addition, categories such as ``plastic wrapper / flexible packaging'' and ``plastic bag / plastic sheet'' can overlap visually. Future releases should publish explicit category decision rules and human agreement on ambiguous examples, together with per-category counts at the independent base-image level.

\paragraph{Controlled rather than naturalistic language.}
Questions are generated through controlled templates and reference answers are drawn from a constrained answer bank. This improves consistency and enables paired bilingual comparisons but yields less linguistic diversity than human-authored VQA. The current benchmark contains 2,412 unique question surface forms but only 100 distinct reference answer strings, so models may benefit from repeated answer semantics. This design can also make generic environmental questions partly answerable without an image, which means the small environmental VGG is partly a property of benchmark construction as well as model behavior. Future versions should incorporate independently authored references, greater paraphrastic diversity, and counterfactual or category-specific questions whose correct answer changes with the visual evidence.

\paragraph{Annotation and bilingual validity.}
Benchmark construction is semi-automated rather than independently expert-annotated at the full-dataset level. This creates a potential confound between model multilingual weakness and unnatural or imperfect Bangla wording. Before release as a high-stakes public benchmark, the full test set or a clearly specified stratified subset should be reviewed by native Bangla annotators for question naturalness and semantic equivalence, and by annotators for object category, visual grounding, and scientific correctness.

\paragraph{Cross-task comparability.}
Recognition, visual description, and environmental reasoning differ not only in semantic demands but also in answer length, answer space, evaluation metric, and required domain knowledge. Differences across these task types therefore should not be interpreted as a controlled causal effect of ``reasoning complexity.'' The normalized score is a reporting convenience that places the primary task-specific metrics on a common numerical scale; it does not make the metrics psychometrically equivalent.

\paragraph{Environmental scope.}
Environmental answers intentionally remain general and scientifically conservative. The benchmark does not evaluate polymer chemistry, life-cycle assessment, detailed ecological toxicity, local waste regulations, or contamination-specific risks. Models should therefore not be interpreted as environmental-science experts on the basis of this benchmark.

\paragraph{Model drift.}
Closed VLM APIs may change over time. Reproducibility requires fixed snapshots where possible and detailed recording of model identifiers, inference dates, prompts, decoding parameters, and provider-side changes.

\paragraph{Metric validity.}
BERTScore, BLEU, and ROUGE compare generated responses with textual references; none guarantees factual correctness. Their behavior can also differ across English and Bangla, making the exact BERTScore encoder and configuration important for interpreting language gaps. Similarly, LaBSE similarity may classify two fluent but factually incorrect responses as cross-lingually consistent, and BCS depends on the selected semantic threshold. Human review remains necessary for high-confidence conclusions.

\paragraph{Dependence among augmented samples.}
Augmented images derived from the same base image are statistically dependent. Confidence intervals based on individual image records or QA pairs would therefore be artificially narrow. Our source-group cluster bootstrap explicitly addresses this dependence.

% NeurIPS 2026 requires the official paper checklist in the same PDF.
% Copy checklist.tex from the official 2026 author kit, answer every item, and keep the questions unchanged.
\newpage
\input{checklist.tex}

\end{document}